\scanpage{167}
Les conditions \deleted{\unsure{aussi}} analogues à (i), (ii), (iii)
$(\mathrm{iii}_0)$ \margin{$(\mathrm{iii}_0)\longrightarrow\text{normalisation}$}.

\deleted{Lorsqu'on dispose du \unsure{supplément} de structure
\unsure{exprimé} par un $\Lambda(-1)$,}
On a un iso
\[
B(M(-1),M(-1))\simeq B(M,M)
\]
et il y a lieu de vérifier la condition que $\Psi_{M(-1)}$ correspond à
$\Psi_M$. Ça ne doit pouvoir faire de difficulté.
\deleted{Pour l'expliciter en termes de $\mathcal V$ directement,
il faudrait expliciter comment \unsure{on obtient} $\Lambda(-1)$
en termes d'un objet.}
En fait, utilisant la propriété (ii), on trouve qu'il suffit de vérifier que
la forme canonique $1$ sur
\[
1\in B(\unsure{M|\Lambda}(-1),\Lambda(-1))
\]
est dans $\Psi_{\Lambda(-1)}$.
\deleted{Ça ne saurait} La vérification dans le cas de Géométrie Alg. est triviale.

Ceci permet d'étendre la notion de forme \inserted{\uline{définie}}
positive aux motifs non nécessairement effectifs.
\mspage{166}
